\documentclass[a4paper,12pt]{article}
\usepackage[T2A]{fontenc}
\usepackage[utf8]{inputenc}
\usepackage[russian]{babel}
\usepackage{amsmath,amssymb,mathtools}
\usepackage{icomma}
\usepackage[margin=22mm,top=21mm,bottom=23mm]{geometry}
\usepackage{microtype}
\usepackage{tikz}
\usetikzlibrary{arrows.meta,positioning,shapes.geometric,calc}
\usepackage{tabularx,booktabs,array}
\usepackage{enumitem}
\usepackage{parskip}
\usepackage{fancyhdr}
\usepackage{setspace}
\usepackage{xcolor}
\usepackage[hidelinks]{hyperref}
\definecolor{accent}{HTML}{008080}
\definecolor{darkaccent}{HTML}{005D5D}
\definecolor{textgray}{HTML}{303538}
\definecolor{softgray}{HTML}{626B70}
\definecolor{rulegray}{HTML}{D2DCDD}
\color{textgray}
\setstretch{1.23}
\setlength{\parindent}{0pt}
\setlength{\parskip}{5pt}
\setlist[itemize]{leftmargin=1.65em,itemsep=1.5pt,topsep=2pt}
\setlist[enumerate]{leftmargin=1.7em,itemsep=8pt,topsep=5pt}
\newcommand{\key}[1]{\textcolor{darkaccent}{\textbf{#1}}}
\newcommand{\mysection}[1]{\vspace{5pt}\section*{\textcolor{darkaccent}{#1}}\vspace{-5pt}}
\newcommand{\myexample}[1]{\textbf{Пример.} #1}
\newcommand{\ruleline}{\vspace{3pt}\noindent\textcolor{rulegray}{\rule{\linewidth}{0.45pt}}\vspace{4pt}}
\newcommand{\xbar}{\overline{x}}
\pagestyle{fancy}
\fancyhf{}
\fancyhead[L]{\footnotesize\color{softgray} Описательная статистика}
\fancyhead[R]{\footnotesize\color{softgray} Николь Саркисьянц}
\fancyfoot[C]{\footnotesize\color{softgray}\thepage}
\renewcommand{\headrulewidth}{0.35pt}
\renewcommand{\footrulewidth}{0pt}
\setlength{\headheight}{15pt}
\begin{document}
\thispagestyle{empty}
\begin{center}
\vspace*{27mm}
{\large\color{darkaccent}\MakeUppercase{Математика}}\par
\vspace{24mm}
{\fontsize{24}{29}\selectfont Чек-лист}\par
\vspace{18mm}
{\fontsize{21}{27}\selectfont\color{accent} Описательная статистика}\par
\vspace{6mm}
{\fontsize{15}{22}\selectfont Среднее, медиана, мода, частоты\\ и характеристики разброса}\par
\vspace{19mm}
{\normalsize Подготовка к ЕГЭ по математике}\par
\vfill
{\large Николь Саркисьянц}\par
\vspace{12mm}
{\normalsize 2 ноября 2026 г.}\par
\vspace{17mm}
{\small Лёвин Артём Александрович\\эксклюзивно для Николь Саркисьянц}\par
\vspace{19mm}
\begin{tikzpicture}[x=1cm,y=1cm]
\draw[accent,line width=0.75pt] (0,0) .. controls (1.4,0.57) and (2.8,-0.46) .. (4.3,0) .. controls (5.6,0.52) and (7.1,-0.46) .. (8.6,0);
\fill[darkaccent] (4.3,0) circle (1.05pt);
\end{tikzpicture}
\end{center}
\clearpage

\mysection{1. Что описывает статистика}
Статистика помогает компактно охарактеризовать \key{набор числовых данных}: определить типичное значение, выяснить, какие числа встречаются чаще, и оценить разброс.

\key{Первый шаг --- упорядочить ряд по неубыванию}, сохраняя все повторения.
\[
1,\;3,\;5,\;4,\;10,\;2,\;15
\quad\longrightarrow\quad
1,\;2,\;3,\;4,\;5,\;10,\;15.
\]
Число элементов --- \key{объём выборки} $N$. В примере $N=7$.

\mysection{2. Среднее арифметическое}
Среднее арифметическое --- сумма всех значений, делённая на их количество:
\[
\boxed{\displaystyle \xbar=\frac{x_1+x_2+\cdots+x_N}{N}.}
\]
\myexample{Для ряда $1,2,3,4,5,10,15$:}
\[
\xbar=\frac{1+2+3+4+5+10+15}{7}=\frac{40}{7}\approx5{,}71.
\]
Среднее не обязано совпадать с каким-либо числом ряда, но находится между его минимумом и максимумом.

\mysection{3. Медиана}
\key{Медиана} $Me$ --- центральное значение \emph{упорядоченного} ряда.
\begin{itemize}
\item При \textbf{нечётном} $N$ берём число ровно посередине.
\item При \textbf{чётном} $N$ находим среднее арифметическое двух центральных чисел.
\end{itemize}
\[
\begin{aligned}
1,2,3,\underbrace{4}_{Me},5,10,15&\quad\Rightarrow\quad Me=4;\\[3pt]
1,2,\underbrace{3,4}_{\text{середина}},5,10&\quad\Rightarrow\quad Me=\frac{3+4}{2}=3{,}5.
\end{aligned}
\]
Повторяющиеся числа тоже учитываются как отдельные элементы ряда.

\mysection{4. Размах}
\key{Размах} --- разность наибольшего и наименьшего значений:
\[
\boxed{R=x_{\max}-x_{\min}.}
\qquad
R=15-1=14.
\]
\clearpage
\begin{center}
{\Large\color{darkaccent}\bfseries Алгоритм: как найти медиану}\par
\vspace{13mm}
\begin{tikzpicture}[
 >=Latex,
 line width=.85pt,
 term/.style={draw=darkaccent,ellipse,minimum width=49mm,minimum height=11mm,align=center,font=\normalsize,inner sep=4pt},
 op/.style={draw=darkaccent,rounded corners=3pt,minimum width=52mm,minimum height=17mm,align=center,text width=48mm,font=\small,inner sep=4pt},
 dec/.style={draw=darkaccent,diamond,aspect=2.15,minimum width=55mm,minimum height=18mm,align=center,font=\small,inner sep=1pt},
 link/.style={->,draw=darkaccent}
]
\node[term] (start) at (0,0) {Начало};
\node[op] (order) at (0,-1.85) {Расположите числа\\по неубыванию};
\node[dec] (dec) at (0,-4.10) {$N$ нечётно?};
\node[op] (odd) at (-3.05,-6.85) {Возьмите одно\\центральное число};
\node[op] (even) at (3.05,-6.85) {Найдите среднее\\двух центральных чисел};
\node[term] (fin) at (0,-9.35) {Медиана найдена};
\draw[link] (start)--(order);
\draw[link] (order)--(dec);
\draw[link] (dec.west)-| node[above,pos=.25]{\small да} (odd.north);
\draw[link] (dec.east)-| node[above,pos=.25]{\small нет} (even.north);
\draw[link] (odd.south)|-(fin.west);
\draw[link] (even.south)|-(fin.east);
\end{tikzpicture}
\end{center}
\clearpage

\mysection{5. Мода и повторяющиеся значения}
\key{Мода} $Mo$ --- значение, встречающееся чаще остальных. Например,
\[
1,\;1,\;3,\;3,\;3,\;5,\;5,\;10,\;10.
\]
Здесь $Mo=3$, так как тройка встречается три раза, а остальные значения --- по два раза. Среднее можно вычислять с учётом количества повторов:
\[
\xbar=\frac{2\cdot1+3\cdot3+2\cdot5+2\cdot10}{9}=\frac{41}{9}.
\]
Если максимальную частоту имеют несколько различных значений, мод несколько. Если все значения встречаются одинаково часто, выделенной моды нет (по используемой здесь школьной договорённости).

\mysection{6. Как выброс влияет на среднее и медиану}
Сравните два упорядоченных ряда:
\[
\begin{array}{c|c|c}
\text{Ряд} & \text{Среднее} & \text{Медиана}\\ \hline
2,3,4,5,6 & 4 & 4\\
2,3,4,5,6,100 & 20 & 4{,}5
\end{array}
\]
Добавление очень большого числа $100$ заметно увеличило среднее, но мало изменило медиану. Такие резко выделяющиеся значения называют \key{выбросами}. Медиана обычно устойчивее к выбросам.

\clearpage
\mysection{7. Абсолютная и относительная частоты}
Для каждого различного значения $x_i$ можно подсчитать:
\begin{itemize}
\item \key{абсолютную частоту} $n_i$ --- сколько раз оно встречается;
\item \key{относительную частоту} $p_i=n_i/N$ (или $100p_i\%$).
\end{itemize}
\[
\boxed{\displaystyle N=\sum_i n_i,\qquad p_i=\frac{n_i}{N},\qquad \sum_i p_i=1.}
\]
Пример: единица встретилась $3$ раза, двойка --- $2$, тройка --- $5$, четвёрка --- $1$. Всего $N=11$.
\begin{center}
\renewcommand{\arraystretch}{1.27}
\begin{tabular}{@{}c c c c@{}}
\toprule
Значение $x_i$ & Частота $n_i$ & Доля $p_i$ & В процентах\\
\midrule
$1$ & $3$ & $3/11$ & $27{,}27\%$\\
$2$ & $2$ & $2/11$ & $18{,}18\%$\\
$3$ & $5$ & $5/11$ & $45{,}45\%$\\
$4$ & $1$ & $1/11$ & $9{,}09\%$\\
\midrule
Итого & $11$ & $1$ & $\approx100\%$\\
\bottomrule
\end{tabular}
\end{center}
\textbf{Важно:} при округлении процентных долей их сумма может немного отличаться от $100\%$.

\mysection{8. Дисперсия --- измеряем разброс}
\key{Дисперсия} $D$ --- среднее арифметическое квадратов отклонений чисел от среднего:
\[
\boxed{\displaystyle D=\frac{(x_1-\xbar)^2+\cdots+(x_N-\xbar)^2}{N}.}
\]
Если данные сгруппированы по частотам, пользуемся формулой
\[
\boxed{\displaystyle \xbar=\frac{\sum_i n_ix_i}{N},\qquad D=\frac{\sum_i n_i(x_i-\xbar)^2}{N}.}
\]
Ту же дисперсию можно вычислять короче:
\[
\boxed{\displaystyle D=\frac{\sum_i n_ix_i^2}{N}-\xbar^{\,2}.}
\]
\textbf{Пошаговый пример.} Возьмём таблицу частот из предыдущего раздела ($N=11$).
\[
\xbar=\frac{3\cdot1+2\cdot2+5\cdot3+1\cdot4}{11}=\frac{26}{11}\approx2{,}36.
\]
Вычислим сумму квадратов с учётом частот:
\[
\sum_i n_ix_i^2=3\cdot1^2+2\cdot2^2+5\cdot3^2+1\cdot4^2=72.
\]
Тогда
\[
D=\frac{72}{11}-\left(\frac{26}{11}\right)^2
=\frac{116}{121}\approx0{,}959.
\]
\key{Проверка смысла:} $D\ge0$. Чем сильнее значения отклоняются от среднего, тем больше обычно дисперсия. Если все числа одинаковы, $D=0$.

\mysection{9. Стандартное отклонение}
\key{Стандартное отклонение} $\sigma$ --- арифметический квадратный корень из дисперсии:
\[
\boxed{\displaystyle \sigma=\sqrt{D}.}
\]
Для предыдущего примера
\[
\sigma=\sqrt{\frac{116}{121}}=\frac{2\sqrt{29}}{11}\approx0{,}98.
\]
Стандартное отклонение выражено в тех же единицах, что и исходные данные; дисперсия --- в их квадратах.

\clearpage
\mysection{10. Интервал «среднее $\pm$ стандартное отклонение»}
Его границы вычисляются так:
\[
\boxed{[\xbar-\sigma\,;\,\xbar+\sigma].}
\]
Для данных из таблицы частот получаем приблизительно
\[
[\,2{,}36-0{,}98\,;\,2{,}36+0{,}98\,]
\approx [\,1{,}38\,;\,3{,}34\,].
\]
В нашем примере внутрь попадают значения $2$ и $3$. Их суммарная относительная частота равна
\[
\frac{2+5}{11}=\frac7{11}\approx63{,}6\%.
\]
\textbf{Точная интерпретация.} Интервал показывает расстояние в одно стандартное отклонение от среднего. Он \emph{не гарантирует}, что новое случайно выбранное значение попадёт в эти границы. Для оценки конкретных вероятностей нужно знать распределение или частоты.

\mysection{11. Универсальный маршрут решения}
\begin{enumerate}[label=\textcolor{darkaccent}{\arabic*.},itemsep=3pt]
\item Упорядочьте значения, сохранив повторения; найдите $N$.
\item Определите медиану, моду (если есть), минимум, максимум и размах.
\item При повторах составьте таблицу $x_i$, $n_i$, $p_i$ и проверьте $\sum n_i=N$.
\item Вычислите среднее по всем значениям или по частотам.
\item Найдите дисперсию: сначала среднее, затем квадраты отклонений или формулу через $x_i^2$.
\item Извлеките квадратный корень, если требуется стандартное отклонение.
\item При необходимости найдите интервал $[\xbar-\sigma\,;\,\xbar+\sigma]$.
\end{enumerate}
\par\smallskip\hrule\smallskip
\textbf{Чего избегать:} медиана без сортировки; потеря повторов; деление на количество различных значений вместо $N$; пропуск частот при вычислении дисперсии; ошибочная трактовка интервала $\xbar\pm\sigma$ как границ всех значений.

\mysection{Памятка для самопроверки}
\begin{itemize}
\item[$\square$] Различаю среднее, медиану, моду и размах.
\item[$\square$] Учитываю повторы, нахожу $N$, абсолютные и относительные частоты.
\item[$\square$] Вычисляю дисперсию, стандартное отклонение и интервал $\xbar\pm\sigma$.
\end{itemize}

\clearpage
\mysection{Тренировочный блок --- Путь А}
\textit{10 заданий. Выполняйте их по порядку: от базовых вычислений к полному описанию выборки. При необходимости округляйте до сотых.}
\begin{enumerate}[label=\textbf{\arabic*.},itemsep=12pt]
\item Упорядочьте числа $8,\;2,\;6,\;4$. Найдите медиану и размах.
\item Для чисел $1,\;3,\;5,\;7,\;9$ вычислите среднее арифметическое и медиану.
\item Для ряда $5,\;5,\;5,\;10,\;10,\;15$ найдите моду, медиану, среднее и размах.
\item Даны числа $-4,\;-1,\;-1,\;2,\;4,\;9$. Найдите среднее, медиану, моду и размах.
\item Среднее арифметическое пяти чисел $5,\;11,\;x,\;15,\;19$ равно $12$. Найдите $x$.
\item Даны числа $3,\;4,\;5,\;6,\;7$. К ряду добавили число $105$. Найдите среднее и медиану \emph{до и после} добавления.
\item Дан ряд $4,\;2,\;4,\;6,\;2,\;4$. Составьте таблицу абсолютных и относительных частот; найдите среднее и моду.
\item Для ряда $2,\;4,\;6,\;8$ вычислите дисперсию, стандартное отклонение и интервал $[\xbar-\sigma\,;\,\xbar+\sigma]$.
\item В выборке число $1$ встречается $2$ раза, число $3$ --- $3$ раза, число $5$ --- $1$ раз. По таблице частот вычислите среднее, дисперсию и стандартное отклонение.
\item В выборке числа $1,3,5,7$ имеют абсолютные частоты $3,2,3,2$ соответственно. Найдите $N$, моду (или моды), медиану, среднее, размах, дисперсию, стандартное отклонение и интервал $[\xbar-\sigma\,;\,\xbar+\sigma]$.
\end{enumerate}

\clearpage
\mysection{Ответы к тренировочному блоку}
\renewcommand{\arraystretch}{1.38}
\small
\begin{tabularx}{\linewidth}{@{}>{\bfseries}p{0.9cm}X@{}}
\toprule
№ & Ответ\\
\midrule
1 & $2,4,6,8$; $Me=5$; $R=6$.\\
2 & $\xbar=5$; $Me=5$.\\
3 & $Mo=5$; $Me=7{,}5$; $\xbar=25/3\approx8{,}33$; $R=10$.\\
4 & $\xbar=1{,}5$; $Me=0{,}5$; $Mo=-1$; $R=13$.\\
5 & $x=10$.\\
6 & До добавления: $\xbar=5$, $Me=5$; после: $\xbar=65/3\approx21{,}67$, $Me=5{,}5$.\\
7 & Для значений $2,4,6$: $n_i=2,3,1$; $p_i=1/3,1/2,1/6$; $\xbar=11/3$; $Mo=4$.\\
8 & $D=5$; $\sigma=\sqrt5\approx2{,}24$; интервал $[5-\sqrt5\,;\,5+\sqrt5]\approx[2{,}76\,;\,7{,}24]$.\\
9 & $\xbar=8/3$; $D=17/9$; $\sigma=\sqrt{17}/3\approx1{,}37$.\\
10 & $N=10$; $Mo\in\{1;5\}$; $Me=4$; $\xbar=3{,}8$; $R=6$; $D=4{,}96$; $\sigma=2\sqrt{31}/5\approx2{,}23$; интервал $[\,1{,}57\,;\,6{,}03\,]$.\\
\bottomrule
\end{tabularx}
\normalsize
\vspace{14pt}
\mysection{Короткая памятка}
\begin{tabularx}{\linewidth}{@{}p{.36\linewidth}X@{}}
\toprule
Если нужно найти\ldots & Используйте\ldots\\
\midrule
типичное значение & среднее, медиану, моду (в зависимости от данных);\\
крайние значения и разброс & минимум, максимум, размах;\\
долю наблюдений & относительную частоту $n_i/N$;\\
разброс вокруг среднего & дисперсию $D$ и стандартное отклонение $\sigma$;\\
интервал в одно $\sigma$ & границы $\xbar-\sigma$ и $\xbar+\sigma$.\\
\bottomrule
\end{tabularx}
\vfill
\begin{center}\footnotesize\color{softgray} Лёвин Артём Александрович эксклюзивно для Николь Саркисьянц\end{center}
\end{document}
